% answers_short.tex —— 2003 年·北京卷·文科 · 简短答案速查
% 由 tools/make_short_answers.py 自 answers.tex 的 \gksolhead/\gkanswer 提取，
% 并与 2003/2003_数学真题_北京文_答案要点.md 逐题比对后生成，勿手改（改源 MD 后重跑）。

\begin{gknotice}
\textbf{使用说明}：本节只给最终答案便于快速核对，完整推导见解析版。答案为编写者独立推导，\textbf{非官方标准答案}。
\end{gknotice}

\gkpart{选择题}
\begin{center}\begin{tabular}{ccccccccccc}
\toprule
题号 & 1 & 2 & 3 & 4 & 5 & 6 & 7 & 8 & 9 & 10 \\
\midrule
答案 & A & D & A & B & D & B & C & B & B & C \\
\bottomrule
\end{tabular}\end{center}

\gkpart{填空题}
\begin{itemize}[leftmargin=2.6em]
  \item[\textbf{11.}] 3
  \item[\textbf{12.}] \(f(x)\) 与 \(g(x)\)
  \item[\textbf{13.}] \(y^2=-36(x-4)\)
  \item[\textbf{14.}] \(\dfrac4{\pi+4}\)
\end{itemize}

\gkpart{解答题}
\begin{itemize}[leftmargin=2.6em]
  \item[\textbf{15.}] (1) 最小正周期为 \(\pi\)；(2) 最大值为 \(\sqrt2\)，最小值为 \(-\sqrt2\)
  \item[\textbf{16.}] (1) \(a_n=2n\)；(2) \(\displaystyle\sum_{k=1}^{n}b_k=\frac{3+(2n-1)3^{n+1}}2\)
  \item[\textbf{17.}] (1) \(A_1D\perp B_1C_1\)；(2) 点 \(D\) 到平面 \(ACC_1\) 的距离为 \(\dfrac{\sqrt3a}{4}\)；(3) \(A_1B\parallel\) 平面 \(ADC_1\)
  \item[\textbf{18.}] (1) U/PENDING；按图示顶点对应时，椭圆为 \(\dfrac{x^2}{25}+\dfrac{y^2}{9}=1\)，准线为 \(x=\pm\dfrac{25}{4}\)；(2) U/PENDING；按图示 \(A=(5,0)\) 对应题干右顶点时，交点 \(M\) 在题给双曲线上
  \item[\textbf{19.}] (1) \(P=(0,4)\) km；(2) \(P=(0,\dfrac{119}{24})\) km
  \item[\textbf{20.}] (1) \(x-1\le f(x)\le1-x\)；(2) 给定函数满足题设条件；(3) 按原页等式逐字理解，不存在满足条件的函数
\end{itemize}
